% Lean module: MonadFormal.Repair
\section{Obtainable repair by measurement}\label{sec:obtainable-repair}

The non-recovery result closes one route: the retained quotient alone
cannot reconstruct a distinction it has identified away. The
remaining question is whether some still-accessible source can supply
the missing distinction, and at what cost.

Suppose the actual information state restricts the latent state to
$S\subseteq X$. An allowed measurement $m_j:X\to O$ costs
$w_j\in\mathbb N$. Measurements are exact, nondestructive reads of
a source the system can still access. A required answer is $r:X\to V$.
An adaptive tree either returns a value or chooses $j$ and branches on
$m_j(x)$. Correctness means that every $x\in S$ receives answer $r(x)$.
Its budget is a worst-case bound along every branch, not an expected cost.
\leandef{MonadFormal.Repair.Tree,MonadFormal.Repair.evaluate,MonadFormal.Repair.Within,MonadFormal.Repair.Correct,MonadFormal.Repair.Feasible}

Write $\mathsf F(S,b)$ when such a correct tree exists within budget $b$,
and $S_{j,o}=S\cap m_j^{-1}(\{o\})$.
\begin{theorem}[Adaptive repair feasibility]\label{thm:repair-feasibility}
The exact recurrence is
\begin{align*}
 \mathsf F(S,b)\iff{}&
   (\exists v\in V:\ \forall x\in S,\ r(x)=v)\\
 &\quad\text{or}\quad
   (\exists j:\ w_j\le b\ \text{and}\
       \forall o\in O,\ \mathsf F(S_{j,o},b-w_j)).
\end{align*}
\end{theorem}
\begin{proof}
A correct tree either returns a constant answer, giving the first
case, or begins with a query $j$. In the latter case the budget pays
$w_j$, and each child must be correct on the compatible fiber
$S_{j,o}$ with the residual budget. Conversely, a constant answer
constructs a leaf. In the second case, choose a witnessing tree for
each outcome and attach it beneath a query node. Its correctness and
budget follow from the corresponding properties of its children.
\leanproof{MonadFormal.Repair.correct_query_iff,MonadFormal.Repair.feasible_iff}
\end{proof}
This is an existence recurrence. Section~\ref{sec:optimal-repair} proves
minimum-cost attainment and a cardinality-founded recursion after pruning
nonsplitting queries. The present tree budget includes unreachable branches;
with a nonempty answer type those branches can receive zero-cost leaves.

\begin{theorem}[Adaptive non-recovery]\label{thm:adaptive-no-recovery}
If $x,y\in S$ agree on every available measurement but
$r(x)\ne r(y)$, no adaptive measurement tree is correct on $S$.
\end{theorem}
\begin{proof}
Induct on the tree. A leaf gives both states the same answer. At a
query node the two measurement values agree, so both states enter the
same child; the induction hypothesis gives equal answers there.
Correctness would therefore imply $r(x)=r(y)$, a contradiction.
\leanproof{MonadFormal.Repair.indistinguishable_evaluation,MonadFormal.Repair.no_repair_from_identical_measurements}
\end{proof}

\begin{theorem}[Zero-cost repair and free information]\label{thm:zero-cost-repair}
If every query costs strictly more than zero, $\mathsf F(S,0)$ holds
exactly when a constant answer is correct on $S$. Without positivity,
a zero-budget tree still cannot distinguish states that agree on every
zero-cost measurement.
\end{theorem}
\begin{proof}
A positive-cost query cannot be paid for at budget zero, so a feasible
tree must be a leaf; conversely, a correct leaf costs zero. For the
second statement, induct on a zero-budget tree. Every queried index
has cost zero. Equal free-measurement values select the same child,
which is still a zero-budget tree.
\leanproof{MonadFormal.Repair.zero_budget_iff_constant,MonadFormal.Repair.zero_budget_respects_free_fibers}
\end{proof}

\begin{example}[A free query defeats a storage-only criterion]
\label{ex:free-query}
Store only a constant from $X=\{0,1\}$ but retain a zero-cost
measurement of the original bit. One query recovers the bit with zero
cost, although the identity probe does not factor through the stored
constant.
\end{example}
\begin{proof}
Query the bit and return the observed value. Both leaves cost zero,
and the answer equals the state. A decoder of the constant storage
would assign the same answer to zero and one, which is impossible.
\leanproof{MonadFormal.Repair.free_bit_costs_zero,MonadFormal.Repair.free_bit_recovers,MonadFormal.Repair.free_bit_not_in_constant_storage}
\end{proof}
Consequently ``zero repair cost'' is relative to the full free-access
structure, not automatically to the stored quotient alone.

% Lean module: MonadFormal.Repair.Optimal
\subsection{Minimum cost, productive queries, and termination}
\label{sec:optimal-repair}

Keep the access model of Section~\ref{sec:obtainable-repair}; in particular,
queries read an actually available source without changing it. Assume the
answer set $V$ is nonempty. Costs are nonnegative integers; free queries
are allowed. Define
\[
 \mathsf V(S)=\inf\{b\in\mathbb N:\mathsf F(S,b)\}
 \quad\text{in }\mathbb N\cup\{\infty\},
 \qquad \inf\varnothing=\infty.
\]
A query is \emph{productive on $S$} when it separates two compatible states:
\[
 J_S=\{j:\exists x,y\in S,\ m_j(x)\ne m_j(y)\}.
\]
For a finite outcome set $O$, a tree has the explicit worst-case cost
\[
 c(\operatorname{answer}(v))=0,\qquad
 c(\operatorname{query}(j,(T_o)_{o\in O}))
   =w_j+\max_{o\in O}c(T_o),
 \qquad \max\varnothing=0.
\]
This counts every syntactic branch. Unreachable branches can receive
arbitrary zero-cost answers because $V$ is nonempty.
\leandef{MonadFormal.Repair.optimalCost,MonadFormal.Repair.Splits,MonadFormal.Repair.worstCost}

\begin{theorem}[Attained minimum repair cost]\label{thm:minimum-repair}
For every finite budget $b$,
\[
 \mathsf V(S)\le b\iff\mathsf F(S,b).
\]
Moreover, $\mathsf V(S)=\infty$ exactly when no finite budget permits
repair. If $O$ is finite and $\mathsf V(S)<\infty$, there is a correct
tree $T$ with $c(T)=\mathsf V(S)$.
\end{theorem}
\begin{proof}
Feasibility is upward closed in the budget: an induction on the tree
preserves its root allowance and increases every residual allowance.
A nonempty set of natural-number budgets has a least member, so the
infimum is a feasible budget whenever it is finite. This proves the
displayed equivalence and the impossibility statement.

For finite $O$, induction on the tree shows that its recursive budget
condition is equivalent to $c(T)\le b$: a query can be paid for precisely
when its cost plus the maximum child cost is at most $b$. Choose a tree
feasible at the least budget. Its cost cannot be larger than that budget;
if it were smaller, the same tree would contradict minimality.
\leanproof{MonadFormal.Repair.within_mono,MonadFormal.Repair.feasible_mono,MonadFormal.Repair.within_iff_worstCost_le,MonadFormal.Repair.optimalCost_le_iff,MonadFormal.Repair.optimalCost_attained,MonadFormal.Repair.optimalCost_eq_top_iff,MonadFormal.Repair.optimal_tree}
\end{proof}

\begin{theorem}[Productive-query Bellman equation]\label{thm:repair-bellman}
If $r$ is constant on $S$, then $\mathsf V(S)=0$; this includes
$S=\varnothing$. Otherwise,
\[
 \mathsf V(S)=\inf_{j\in J_S}
   \left[w_j+\sup_{o\in O}\mathsf V(S_{j,o})\right].
\]
For finite query and outcome sets, these are a minimum and maximum,
with the empty-set conventions above. No positive-cost assumption is
required. In particular, a nonconstant target can still have value zero
when productive free measurements are available.
\end{theorem}
\begin{proof}
A constant answer gives a zero-cost leaf. Otherwise $S$ is nonempty.
Consider a feasible tree. If its root query does not split $S$, every
compatible state selects the same child. Delete that root and retain
this child: correctness on $S$ is unchanged, and the child's allowance
can be enlarged to the original budget. Induction on the original tree
eventually exposes a productive root, since a correct constant leaf
would contradict nonconstancy. Conversely, any productive root with
feasible children gives a feasible tree. Thus
\[
 \mathsf F(S,b)\iff
 \exists j\in J_S:\ w_j\le b\ \text{and}\
   \forall o,\ \mathsf F(S_{j,o},b-w_j).
\]
Replace child feasibility by the corresponding bound on $\mathsf V$.
For extended natural numbers, the bound on a query's cost plus its
supremum child value is equivalent to these residual-budget bounds.
An infimum in $\mathbb N\cup\{\infty\}$ has a witness whenever it is
at most a finite $b$. Consequently both sides of the Bellman equation
have exactly the same finite upper bounds, which determines their value.
\leanproof{MonadFormal.Repair.feasible_split_iff,MonadFormal.Repair.inf_le_nat_iff,MonadFormal.Repair.nat_add_le_nat_iff,MonadFormal.Repair.eq_of_same_nat_bounds,MonadFormal.Repair.optimalCost_bellman,MonadFormal.Repair.optimalCost_constant,MonadFormal.Repair.optimalCost_empty}
\end{proof}

\begin{theorem}[Cardinality-founded repair recursion]\label{thm:repair-termination}
For finite $S$, every branch of a productive query satisfies
$|S_{j,o}|<|S|$. The recursion that returns zero on constant-target sets
and otherwise applies the productive-query Bellman equation is therefore
well founded and returns $\mathsf V(S)$.
\end{theorem}
\begin{proof}
Choose $x,y\in S$ with distinct query outcomes. For any proposed outcome
$o$, at least one of these states is absent from $S_{j,o}$. That branch
is a proper subset of the finite set $S$, hence has smaller cardinality.
Define the recursion using this decreasing measure. Strong induction on
$|S|$ identifies its base case with the constant-answer value and its
recursive case with the Bellman equation.
\leanproof{MonadFormal.Repair.split_branch_card_lt,MonadFormal.Repair.finiteCost_eq_optimal}
\end{proof}
\leandef{MonadFormal.Repair.finiteCost}
For finite $J$ and $O$, this is a finite recursive optimization problem.
Its constant-target test does not require searching an infinite answer
space: it is equivalent to pairwise agreement of $r(x)$ on $S$.
\begin{proof}
A constant target gives pairwise agreement. Conversely, if $S$ is
nonempty, take the target value at any member of $S$; pairwise agreement
makes it a constant answer. If $S$ is empty, choose any member of $V$.
\leanproof{MonadFormal.Repair.constant_iff_pairwise}
\end{proof}
The recurrence establishes neither a low computational complexity bound
nor repair from unavailable measurements. It solves the declared
exact, nondestructive, worst-case query problem.

% Lean module: MonadFormal.Repair.Supplement
\subsection{Exact supplements and the timing of relevance}
\label{sec:supplement-coding}

The next access model is different from measurement repair. An encoder
still sees $x\in X$, while the decoder receives only retained side
information $q(x)\in Q$ and a supplementary message $e(x)\in A$.
The required answer is $r(x)\in V$. An exact supplement consists of
maps $e:X\to A$ and $d:Q\times A\to V$ satisfying
\[
 d(q(x),e(x))=r(x)\qquad(x\in X).
\]
Assume $X$ and $Q$ are finite and $V$ is nonempty. For each retained
value $y$, let
\[
 R_y=r(q^{-1}(\{y\})),\qquad
 M=\max_{y\in Q}|R_y|,
 \qquad \max\varnothing=0.
\]
\leandef{MonadFormal.Repair.HasSupplement,MonadFormal.Repair.fiberTargets,MonadFormal.Repair.supplementSize}

\begin{theorem}[Minimum exact supplement alphabet]\label{thm:supplement-alphabet}
An exact supplement with alphabet $\{0,\ldots,k-1\}$ exists if and
only if $M\le k$. In particular, the minimum alphabet size is $M$.
For fixed-length binary messages, the minimum length is
\[
 b_{\min}=\min\{b\in\mathbb N:M\le 2^b\}.
\]
This is the integer ceiling logarithm, with value zero also when $M=0$.
\end{theorem}
\begin{proof}
Fix a retained value $y$. Every target value in $R_y$ must be obtainable
by applying $d(y,\cdot)$ to a message. Therefore $|R_y|\le |A|$,
which gives the lower bound for every fiber and hence for their maximum.

Conversely, if $|R_y|\le k$, enumerate the distinct target values in
each fiber using distinct indices in $\{0,\ldots,k-1\}$. The encoder
knows $q(x)$ and $r(x)$, so it sends the index of $r(x)$ in that fiber's
enumeration. The decoder uses $y$ to select the enumeration and recover
the target. Give unused indices arbitrary answers in the nonempty set
$V$. Distinct fibers may reuse the same message indices because the
decoder already knows $y$. This constructs an exact supplement of
size $M$ as well as one of any larger size.

There are $2^b$ binary words of length $b$. Relabeling an alphabet
preserves the coding equation, so a length-$b$ code exists exactly
when $M\le2^b$. The least such natural number is the integer ceiling
logarithm; the corresponding code exists, not merely a limiting code.
\leanproof{MonadFormal.Repair.supplement_lower_bound,MonadFormal.Repair.supplement_upper_bound,MonadFormal.Repair.supplement_iff,MonadFormal.Repair.supplementSize_le_iff,MonadFormal.Repair.supplement_minimum,MonadFormal.Repair.supplement_relabel,MonadFormal.Repair.supplement_bits_iff,MonadFormal.Repair.supplement_binary_iff,MonadFormal.Repair.supplement_minimum_bits}
\end{proof}
This is an exact, one-shot, worst-case alphabet result, not an entropy
or asymptotic average-rate assertion. It assumes an encoder with access
to the original state. It supplies no method for reconstructing the
supplement from an already insufficient $q(x)$ alone, and it does not
charge computation or measurement costs to the encoder.

\begin{theorem}[A question known before versus after compression]
\label{thm:query-timing}
Let $X=\{0,1\}^n$. If a coordinate $i$ is known to the encoder before
compression, one bit suffices to answer that coordinate. In contrast,
suppose one common encoding $e:X\to A$ must support exact answers to
every coordinate selected after compression, with no other retained
information or later measurement. Then $e$ is injective and
$|A|\ge2^n$. A binary fixed-length encoding therefore requires at
least $n$ bits, and storing all $n$ bits attains the bound.
\end{theorem}
\begin{proof}
For a known coordinate, send $x_i$ and return the received bit. For a
question chosen later, suppose $e(x)=e(y)$. The same decoder, applied
to either message and any coordinate $i$, must return both $x_i$ and
$y_i$. Thus all coordinates agree and $x=y$. The encoding is injective,
so its alphabet has at least $2^n$ elements. A $b$-bit alphabet has
$2^b$ elements, hence $b\ge n$. The identity encoding stores the vector
and answers every coordinate, establishing attainability.
\leanproof{MonadFormal.Repair.known_query_one_bit,MonadFormal.Repair.all_queries_force_injective,MonadFormal.Repair.all_queries_alphabet_bound,MonadFormal.Repair.all_queries_need_n_bits,MonadFormal.Repair.storing_n_bits_answers_all}
\end{proof}
The distinction is a change in the admitted question family and its
timing, not the creation of a previously nonexistent coordinate. It
provides a precise baseline for anticipatory compression without a
novelty claim or an assumption of access to discarded state.

